\setcounter{section}{24}

  \zwischenueberschrift{Soal latihan}

\inputaufgabe
{}
{

Misalkan $X$ suatu
\definitionsverweis {manifold diferensiabel}{}{}
dengan
\definitionsverweis {basis topologi terhitung}{}{}
dan misalkan

\mathdisp {0 \, \stackrel{  }{\longrightarrow} \,  U   \, \stackrel{  }{ \longrightarrow}   \,  V   \, \stackrel{  }{\longrightarrow} \,  W \,\stackrel{  } {\longrightarrow}  \, 0} {  }

suatu
\definitionsverweis {barisan eksak pendek}{}{}
dari
\definitionsverweis {bundel-bundel vektor diferensiabel}{}{}
di atas $X$. Dengan menggunakan
[[Mannigfaltigkeit/Abzählbare Basis/Überdeckung/Partition/Fakt|Teorema 22.10]],
tunjukkan bahwa barisan tersebut memiliki suatu pemisahan, yaitu suatu isomorfisme bundel vektor

\mavergleichskette
{\vergleichskette
{ V
}
{ \cong }{ U \oplus W
}
{  }{
}
{    }{
}
{   }{
}
}
{}{}{,}
yang kompatibel dengan homomorfisme-homomorfisme bundel.

}
{} {}

\inputaufgabe
{}
{

Misalkan

\mavergleichskette
{\vergleichskette
{ W
}
{ \subseteq }{ \R^n
}
{  }{
}
{    }{
}
{   }{
}
}
{}{}{}
\definitionsverweis {terbuka}{}{,}
\maabb {h} { W } { \R
} {}
suatu
\definitionsverweis {fungsi terdiferensialkan secara kontinu}{}{}
dan

\mavergleichskette
{\vergleichskette
{ Y
}
{ = }{ h^{-1}(c)
}
{  }{
}
{    }{
}
{   }{
}
}
{}{}{}
merupakan
\definitionsverweis {serat}{}{}
di atas

\mavergleichskette
{\vergleichskette
{ c
}
{ \in }{ \R
}
{  }{
}
{    }{
}
{   }{
}
}
{}{}{,}
dan anggap $h$
\definitionsverweis {reguler}{}{}
di setiap titik $Y$. Realisasikan bundel tangen pertama dan kedua
\definitionsverweis {bundel tangen}{}{}
dari $Y$ sebagai
\definitionsverweis {submanifold tertutup}{}{}
di $W \times \R^n$ dan, secara berturut-turut, di
\mathl{W \times \R^n \times \R^n \times \R^n}{} dalam pengertian
[[Differenzierbare Hyperfläche/Zweites Tangentialbündel/Implizite Beschreibung/Bemerkung|Catatan 24.4]]
bersama dengan proyeksi bundel
\maabbeledisp {\pi} {TY} {Y
} {(P,u)} {P
} {,}
dan
\maabbeledisp {q} {TTY} {TY
} {(P,u,v,w)} {(P,u)
} {,}
Tunjukkan pernyataan-pernyataan berikut.
\aufzaehlungdrei{Suatu
\definitionsverweis {kurva diferensiabel}{}{}
\maabbeledisp {\epsilon} { I } { TY
} { t } { (P(t), u(t))
} {,}
pada interval terbuka

\mavergleichskette
{\vergleichskette
{ I
}
{ \subseteq }{ \R
}
{  }{
}
{    }{
}
{   }{
}
}
{}{}{}
dengan

\mavergleichskette
{\vergleichskette
{ 0
}
{ \in }{ I
}
{  }{
}
{    }{
}
{   }{
}
}
{}{}{}
mendefinisikan
\definitionsverweis {vektor singgung}{}{}

\mavergleichskette
{\vergleichskette
{ [\epsilon]
}
{ = }{ (P(0), u(0), P'(0), u'(0))
}
{ \in }{ TTY
}
{    }{
}
{   }{
}
}
{}{}{.}
}{Di bawah
\definitionsverweis {pemetaan tangen}{}{}
\maabbdisp {T_{(P,u)} (\pi)} {T_{ (P,u)} TY} { T_PY
} {}
yang berkaitan dengan $\pi$, kelas
\mathl{[\epsilon]}{} dari (1) dipetakan ke

\mavergleichskettedisp
{\vergleichskette
{ [ \pi \circ \epsilon ]
}
{ =} { [P'(0)]
}
{ } {
}
{ } {
}
{ } {
}
}
{}{}{}
.
}{Di bawah
\definitionsverweis {pemetaan tangen}{}{}
\maabb {T(\pi)} {TTY} {TY
} {}
yang berkaitan dengan $\pi$, elemen
\mathl{(P,u,v,w)}{} dipetakan ke
\mathl{(P,v)}{}.
}

}
{} {}

\inputaufgabe
{}
{

Untuk lingkaran satuan

\mavergleichskettedisp
{\vergleichskette
{ Y
}
{ =} { \left\{ (x_1,x_2)  \mid   x_1^2 + x_2^2 = 1         \right\}
}
{ \subset} { \R^2
}
{ } {
}
{ } {
}
}
{}{}{}
tentukan deskripsi implisit
\definitionsverweis {bundel tangen}{}{}
$TY$ dan bundel tangen kedua $TTY$ dalam pengertian
[[Differenzierbare Hyperfläche/Zweites Tangentialbündel/Induzierte Bündelhomomorphismen/Bemerkung|Catatan 24.5]].

}
{} {}

\inputaufgabe
{}
{

Misalkan

\mavergleichskette
{\vergleichskette
{ W
}
{ \subseteq }{ \R^n
}
{  }{
}
{    }{
}
{   }{
}
}
{}{}{}
\definitionsverweis {terbuka}{}{,}
\maabb {\varphi} { W } { \R^k
} {}
suatu pemetaan
$C^2$-\definitionsverweis {terdiferensialkan secara kontinu}{}{},

\mavergleichskette
{\vergleichskette
{ c
}
{ \in }{ \R^k
}
{  }{
}
{    }{
}
{   }{
}
}
{}{}{}
dan

\mavergleichskette
{\vergleichskette
{ Y
}
{ = }{ \varphi^{-1} (c)
}
{  }{
}
{    }{
}
{   }{
}
}
{}{}{}
merupakan
\definitionsverweis {serat}{}{} di atas $c$. Anggap serat ini
\definitionsverweis {reguler}{}{}
di setiap titiknya. Kembangkan suatu deskripsi implisit
\zusatzklammer {yakni dengan persamaan} {} {}
dari
\definitionsverweis {bundel tangen}{}{}
$TY$ dan bundel tangen kedua $TTY$, serupa dengan
[[Differenzierbare Hyperfläche/Zweites Tangentialbündel/Induzierte Bündelhomomorphismen/Bemerkung|Catatan 24.5]].

}
{} {}

\inputaufgabe
{}
{

Misalkan $X$ suatu
$C^2$-\definitionsverweis {manifold diferensiabel}{}{.}
Misalkan

\mavergleichskette
{\vergleichskette
{ I,J
}
{ \subseteq }{ \R
}
{  }{
}
{    }{
}
{   }{
}
}
{}{}{}
\definitionsverweis {interval terbuka}{}{}
dengan

\mavergleichskette
{\vergleichskette
{ 0
}
{ \in }{ I,J
}
{  }{
}
{    }{
}
{   }{
}
}
{}{}{}
dan misalkan
\maabbeledisp {\varphi} {I \times J} {X
} {(s,t)} { \varphi(s,t)
} {,}
suatu pemetaan dua kali
\definitionsverweis {terdiferensialkan secara kontinu}{}{}
dengan

\mavergleichskette
{\vergleichskette
{ \varphi(0,0)
}
{ = }{ P
}
{ \in }{ X
}
{    }{
}
{   }{
}
}
{}{}{.}
\aufzaehlungvier{Tunjukkan bahwa untuk setiap

\mavergleichskette
{\vergleichskette
{ s
}
{ \in }{ I
}
{  }{
}
{    }{
}
{   }{
}
}
{}{}{}
pembatasan

\mathl{\varphi \vert_{ \{s\} \times J }}{} mendefinisikan suatu kurva diferensiabel di $X$.
}{Tunjukkan bahwa (1) mendefinisikan suatu kurva diferensiabel
\maabbeledisp {} { I } { TX
} { s } { [ \varphi \vert_{ \{s\} \times J } ]
} {,}
di mana
\mathl{[ \varphi \vert_{ \{s\} \times J } ]}{} menyatakan vektor singgung kurva tersebut pada titik $(s,0)$.
}{Tunjukkan bahwa (2) menentukan suatu elemen dalam
\mathl{TTY}{}.
}{Tunjukkan bahwa setiap elemen
\mathl{TTY}{} dapat direalisasikan oleh suatu pemetaan $\varphi$ seperti di atas.
}

}
{} {}

\inputaufgabe
{}
{

Misalkan $X$ suatu
$C^2$-\definitionsverweis {manifold diferensiabel}{}{.}
Kita memandang pemetaan-pemetaan yang dua kali
\definitionsverweis {terdiferensialkan secara kontinu}{}{}

\maabbeledisp {\varphi} {I \times J} {X
} {(s,t)} { \varphi(s,t)
} {,}
sebagai elemen dalam
\mathl{TTX}{} dalam arti
[[Differenzierbare_Mannigfaltigkeit/Quader/Abbildung/Zweites_Tangentialbündel/Aufgabe|Soal 24.5]].
Bagaimana realisasi ini berperilaku berdasarkan barisan eksak pendek

\mathdisp {0 \, \stackrel{  }{\longrightarrow} \, p^*TX  \, \stackrel{  }{ \longrightarrow}   \, TTX  \, \stackrel{  }{\longrightarrow} \, p^*TX \,\stackrel{  } {\longrightarrow}  \, 0} { . }

}
{} {}

\inputaufgabe
{}
{

Misalkan
\maabb {} { E } { X
} {}

suatu
\definitionsverweis {bundel vektor diferensiabel}{}{}
di atas suatu
\definitionsverweis {manifold diferensiabel}{}{}
$X$. Tentukan
\definitionsverweis {rank}{}{}
bundel-bundel vektor di atas $E$ dalam barisan eksak pendek

\mathdisp {0 \, \stackrel{  }{\longrightarrow} \,  V   \, \stackrel{  }{ \longrightarrow}   \,  TE  \, \stackrel{  }{\longrightarrow} \, p^* TX \,\stackrel{  } {\longrightarrow}  \, 0} { . }

}
{} {}

\inputaufgabe
{}
{

Tentukan
\definitionsverweis {koneksi}{}{}
pada bundel nol di atas suatu
\definitionsverweis {manifold diferensiabel}{}{}
$X$.

}
{} {}

\inputaufgabe
{}
{

Tentukan
\definitionsverweis {koneksi}{}{}
pada suatu
\definitionsverweis {bundel vektor}{}{}
di atas
\definitionsverweis {manifold}{}{}
satu titik $X$.

}
{} {}

\inputaufgabe
{}
{

Misalkan $X$ suatu
\definitionsverweis {manifold diferensiabel}{}{}
dan $E$ suatu
\definitionsverweis {bundel vektor diferensiabel}{}{}
di atas $X$ yang dilengkapi dengan
\definitionsverweis {koneksi}{}{}
. Tunjukkan bahwa konsep turunan vertikal juga dapat didefinisikan sepanjang suatu pemetaan diferensiabel
\maabbdisp {\psi} { Z } { X
} {}
; bandingkan pula dengan
[[Hyperfläche/Tangentiales Vektorfeld/Längs einer Kurve/Kovariante Ableitung/Definition|Definisi 6.5]].

}
{} {}

  \zwischenueberschrift{Soal untuk dikumpulkan}

\inputaufgabe
{3}
{

Misalkan $X$ suatu
\definitionsverweis {ruang topologis}{}{}
dan misalkan

\mathdisp {0 \, \stackrel{  }{\longrightarrow} \, U  \, \stackrel{  }{ \longrightarrow}   \, V  \, \stackrel{  }{\longrightarrow} \, W \,\stackrel{  } {\longrightarrow}  \, 0} {  }

suatu
\definitionsverweis {barisan eksak pendek}{}{}
dari
\definitionsverweis {bundel vektor}{}{}
di atas $X$. Misalkan
\maabb {\varphi} {Y} {X
} {}
suatu
\definitionsverweis {pemetaan kontinu}{}{}
dari suatu ruang topologis lain $Y$. Tunjukkan bahwa penarikan balik menghasilkan barisan eksak pendek

\mathdisp {0 \, \stackrel{  }{\longrightarrow} \, \varphi^*U  \, \stackrel{  }{ \longrightarrow}   \, \varphi^*V  \, \stackrel{  }{\longrightarrow} \, \varphi^*W \,\stackrel{  } {\longrightarrow}  \, 0} {  }

bundel vektor di atas $Y$.

}
{} {}

\inputaufgabe
{3 (1+2)}
{

Misalkan

\mavergleichskette
{\vergleichskette
{ I
}
{ \subseteq }{ \R
}
{  }{
}
{    }{
}
{   }{
}
}
{}{}{}

\definitionsverweis {interval terbuka}{}{,}
kita tinjau
\definitionsverweis {barisan eksak pendek}{}{}
\zusatzklammer {dari bundel vektor di atas $I \times \R$} {} {}

\mathdisp {0 \, \stackrel{  }{\longrightarrow} \, I \times \R \times \R   \, \stackrel{  }{ \longrightarrow}   \,  I \times \R \times \R \times \R   \, \stackrel{  }{\longrightarrow} \,  I \times \R \times \R  \,\stackrel{  } {\longrightarrow}  \, 0} { , }

di mana pemetaan pertama diberikan oleh
\mathl{(P,u,w) \mapsto (P,u,0,w)}{} dan pemetaan kedua oleh
\mathl{(P,u,v,w) \mapsto (P,u,v)}{}. Misalkan
\maabbdisp {g,h} {I \times \R } {\R
} {}
fungsi-fungsi terdiferensialkan secara kontinu.
\aufzaehlungdrei{Tunjukkan bahwa
\maabbeledisp {} {I \times \R \times \R \times \R  } { I \times \R \times \R
} {(P,u,v,w)} { (P,u,   v g(P,u) +w h(P,u) )
} {,}
mendefinisikan suatu homomorfisme bundel vektor di atas $I \times \R$.
}{Misalkan

\mavergleichskette
{\vergleichskette
{ h
}
{ = }{ 1
}
{  }{
}
{    }{
}
{   }{
}
}
{}{}{}
konstan. Tunjukkan bahwa pemetaan pada (1) mendefinisikan suatu pemisahan untuk bundel di tengah.
}{
}

}
{} {}

\inputaufgabe
{4}
{

Untuk kurva bidang

\mavergleichskettedisp
{\vergleichskette
{ Y
}
{ =} { \left\{ (x_1,x_2)  \mid   x_1^3 -5x_1x_2 + x_2^4 = 1         \right\}
}
{ \subset} { \R^2
}
{ } {
}
{ } {
}
}
{}{}{}
tentukan deskripsi implisit
\definitionsverweis {bundel tangen}{}{}
$TY$ dan bundel tangen kedua $TTY$ dalam pengertian
[[Differenzierbare Hyperfläche/Zweites Tangentialbündel/Induzierte Bündelhomomorphismen/Bemerkung|Catatan 24.5]].

}
{} {}

\inputaufgabe
{5 (3+2)}
{

Kita tinjau permukaan diferensiabel

\mavergleichskettedisp
{\vergleichskette
{ Y
}
{ =} { \left\{ (x_1,x_2,x_3)  \mid   x_1^2+x_2^2 -x_3^2 = 0  , \, (x,y,z) \neq 0       \right\}
}
{ \subseteq} { W
}
{ =} { \R^3 \setminus \{(0,0,0) \}
}
{ \subset} { \R^3
}
}
{}{}{.}
\aufzaehlungzwei {Tentukan deskripsi implisit dari
\definitionsverweis {bundel tangen}{}{}
$TY$ dan bundel tangen kedua $TTY$ dalam pengertian
[[Differenzierbare Hyperfläche/Zweites Tangentialbündel/Induzierte Bündelhomomorphismen/Bemerkung|Catatan 24.5]].
} {Hitung untuk kurva diferensiabel
\maabbeledisp {} { I } { Y
} { t } { \left(  \cos  t   , \,   \sin  t  , \,   1                 \right)
} {}
kurva terkait di bundel tangen dan di bundel tangen kedua.
}

}
{} {}

\inputaufgabe
{3}
{

Misalkan

\mavergleichskette
{\vergleichskette
{ U
}
{ \subseteq }{ \R^n
}
{  }{
}
{    }{
}
{   }{
}
}
{}{}{}
\definitionsverweis {terbuka}{}{}
dan misalkan

\mavergleichskette
{\vergleichskette
{ E
}
{ = }{ U \times \R
}
{  }{
}
{    }{
}
{   }{
}
}
{}{}{,}
di mana fungsi-fungsi
\maabb {f} { U } {\R
} {}
kita pandang sebagai
\definitionsverweis {penampang}{}{}
di $E$. Lengkapi $E$ dengan
\definitionsverweis {koneksi trivial}{}{.} Tunjukkan bahwa

\mavergleichskette
{\vergleichskette
{ \nabla (f)
}
{ = }{ df
}
{  }{
}
{    }{
}
{   }{
}
}
{}{}{}
untuk fungsi-fungsi terdiferensialkan secara kontinu.

}
{} {}

 [[Kategorie:Latexseite]]
